% TOP_PROBLEM: {"id": 30005276, "title": "Density Hypothesis for Riemann-Zeta Zeros", "category_id": 1, "category": {"id": 1, "name": "number_theory", "display_name": "Number Theory", "description": "Properties of integers, prime numbers, Diophantine equations.", "slug": "number-theory", "order_index": 1, "created_at": "2026-07-31T15:26:25.670Z"}, "status": "partially_solved"}
% FIRST_POSTED: 2026-09-27
% !TeX program = lualatex
% Compile twice: lualatex 30005276.tex
% All references and prose are embedded; no BibTeX database or external inputs.
\documentclass[11pt,a4paper]{article}
\usepackage[margin=24mm,headheight=14pt]{geometry}
\usepackage{fontspec}
\setmainfont{Latin Modern Roman}
\setsansfont{Latin Modern Sans}
\setmonofont{Latin Modern Mono}
\usepackage{amsmath,amssymb,mathtools}
\usepackage{unicode-math}
\setmathfont{Latin Modern Math}
\usepackage{microtype}
\usepackage{enumitem,xurl,needspace}
\usepackage[dvipsnames]{xcolor}
\usepackage{hyperref}
\hypersetup{unicode=true,colorlinks=true,linkcolor=MidnightBlue,urlcolor=MidnightBlue,
 pdftitle={500 Mathematical Problem Statements},pdfauthor={Source-annotated compilation},bookmarksnumbered=true}
\usepackage{bookmark}
\usepackage{tocloft}
\setlength{\cftsecnumwidth}{3em}
\cftsetpnumwidth{2.5em}
\cftsetrmarg{4em}
\usepackage{titlesec}
\titleformat{\section}{\Large\bfseries\raggedright}{\thesection}{0.7em}{}
\usepackage{fancyhdr}
\pagestyle{fancy}\fancyhf{}
\fancyhead[L]{\small\sffamily Mathematical problem statements}
\fancyhead[R]{\small Source edition 22}
\fancyfoot[C]{\thepage}
\setlength{\parindent}{0pt}
\setlength{\parskip}{5pt plus 1pt}
\setlength{\emergencystretch}{3em}
\setcounter{tocdepth}{1}
\setcounter{secnumdepth}{1}
\setlist[itemize]{leftmargin=3em,itemsep=5pt,topsep=4pt,parsep=0pt}
\allowdisplaybreaks
\newcommand{\N}{\mathbb{N}}
\newcommand{\Z}{\mathbb{Z}}
\newcommand{\Q}{\mathbb{Q}}
\newcommand{\R}{\mathbb{R}}
\newcommand{\C}{\mathbb{C}}
\newcommand{\F}{\mathbb{F}}
\newcommand{\A}{\mathbb{A}}
\newcommand{\PP}{\mathbb P}
\newcommand{\E}{\mathbb{E}}
\newcommand{\eps}{\varepsilon}
\newcommand{\dd}{\,\mathrm d}
\newcommand{\sref}[2]{\hyperlink{source:#1:#2}{[S#2]}}
\newcommand{\eref}[2]{\hyperlink{extra:#1:#2}{[E#2]}}
% Turn the next switch off for a shorter reading copy. The quotations preserve
% every exactTarget field, including qualifications omitted from a short summary.
\newif\ifshowcatalogue
\showcataloguetrue
\newcommand{\cataloguescope}[1]{\ifshowcatalogue
 \paragraph{Exact scope in the supplied catalogue.}
 \begin{quote}\small #1\end{quote}\fi}
\hypersetup{pdftitle={Problem 30005276 - Zeta zero density hypothesis},pdfauthor={Open Problems Math research notebook}}
\fancyhead[L]{\small\sffamily Open Problems Math \textbar\ Research notebook}
\fancyhead[R]{\small\sffamily 128 \textbar\ 27 September 2026}
\numberwithin{equation}{section}
\begin{document}
\setcounter{section}{0}
% ================================================================
% problemId: 30005276
% Canonical problem ID: 30005276
% exactTarget SHA-256: 3c1a915de699bd704436e9353a8a97c10069b6c0f9276b4c73eed5335c6dbf59
\clearpage
\section*{\texorpdfstring{Density Hypothesis for Riemann-Zeta Zeros}{Density Hypothesis for Riemann-Zeta Zeros}}\label{30005276}
\subsection{Definitions and mathematical statement}
Count zeta zeros with multiplicity by
\[
 N(\sigma,T)=\#\{\rho=\beta+i\gamma:\zeta(\rho)=0,\ 0<\beta<1,
                                  \beta\ge\sigma,\ |\gamma|\le T\}.
\]
For each fixed $1/2\le\sigma\le1$ and $\varepsilon>0$, the density hypothesis predicts
\[
 N(\sigma,T)\ll_{\sigma,\varepsilon}T^{2(1-\sigma)+\varepsilon}.
\]
The bound is uniform in large $T$, with dependence on $\sigma,\varepsilon$ permitted. It limits the number of off-line zeros rather than asserting that they do not exist.
\par\smallskip\textit{Formulation and concept references:} \sref{30005276}{1}.
\cataloguescope{Prove that for every sigma with 1/2 <= sigma <= 1 and every epsilon > 0, the number N(sigma,T) of nontrivial zeros rho = beta + i gamma of the Riemann zeta function with beta >= sigma and |gamma| <= T satisfies N(sigma,T) <= C\_\{sigma,epsilon\} T\textasciicircum{}\{2(1-sigma)+epsilon\} for all sufficiently large T.}
\subsection{Short English statement}
Are zeta zeros to the right of any fixed vertical line as sparse as the density hypothesis predicts?
% BEGIN RESEARCH ADDITION 128
\subsection{Research attempt}

\paragraph{Current frontier and scope.}
The zero count includes multiplicities and both signs of the ordinate.
Guth--Maynard's Theorem 1.2 gives
$N(\sigma,T)\leq T^{15(1-\sigma)/(3+5\sigma)+o(1)}$; combining it with
Ingham's estimate gives the uniform benchmark
$T^{30(1-\sigma)/13+o(1)}$.  The version consulted was revised on
7 April 2026.  Its Section 13 carefully separates zero detection from
large-value estimation.  Neither displayed exponent is the conjectural
$2(1-\sigma)$ throughout the strip. \eref{30005276}{1}
These are the relevant established benchmarks for this attempt, not a
claim that one formula is the strongest available at every $\sigma$.

\paragraph{Attack map.}
A large-values approach needs stronger control of the particular
Dirichlet polynomials produced by zero detection, not of an arbitrary
coefficient sequence without hypotheses.  A mollifier approach tries
to make $\zeta M$ close to $1$ in mean square: every zero forces a
unit-sized defect, which should consume some of that mean square.
A direct zero-repulsion approach would need a spacing theorem stronger
than the elementary local zero count; simple distinctness or multiplicity
bounds would not control the total.  The mollifier route is developed
below using area integrals, which makes the zero-detection implication
particularly transparent and exposes its still-unproved analytic input.

\paragraph{Attempt I: an area-budget zero detector.}
For $T\geq4$, set
\[
 M_T(s)=\sum_{d\leq\lfloor T\rfloor}\frac{\mu(d)}{d^s},\qquad
 g_T(s)=\zeta(s)M_T(s)-1.
\]
For each fixed $1/2<u_0<1$, define the genuine analytic quantity
\begin{equation}\label{30005276:area}
 \mathcal E(u_0,T)=\int_{u_0}^{2}\int_{T-1}^{2T+1}
                       |g_T(u+it)|^2\,dt\,du.
\end{equation}
No Dirichlet-series expansion in the critical strip is assumed in this
definition.  The following is a sufficient, but \emph{unproved}, estimate:
\begin{equation}\label{30005276:missing}
 \forall u_0\in(1/2,1)\ \forall\eta>0:\qquad
 \mathcal E(u_0,T)\ll_{u_0,\eta}T^{2(1-u_0)+\eta}.
\end{equation}
The mollifier length here grows with $T$.  A theorem for a fixed
mollifier with constants depending arbitrarily on its length would not
establish this assertion.

\textbf{Area-detection lemma.}
Estimate \eqref{30005276:missing} implies the density hypothesis in the exact
form stated in this entry.

\emph{Proof.} Fix $1/2<\sigma<1$ and $\varepsilon>0$.  Choose
\[
 0<r<\min\{(\sigma-1/2)/2,1/4,\varepsilon/8\},\qquad u_0=\sigma-r.
\]
The standard local zero bound is $O(\log T)$ zeros, counted with
multiplicity, in any horizontal strip of height one at ordinates of
size $T$; this is explicitly used in \eref{30005276}{1}, Section 13.1.
From the zeros with $\beta\geq\sigma$ and $T\leq\gamma\leq2T$,
greedily select centers whose ordinates are at least $2r$ apart.
Each choice removes only $O(\log T)$ zeros, including repeated zeros.
Thus, if the original count is $Z$, there are at least
$Z/(C\log T)$ selected centers.

For each selected $\rho$, the disk $D(\rho,r)$ lies inside the rectangle
in \eqref{30005276:area}.  The disks are pairwise disjoint up to boundaries.
The function $g_T$ is analytic there, since the pole of $\zeta$ is far
away, and $g_T(\rho)=-1$.  The submean inequality for the subharmonic
function $|g_T|^2$ gives
\[
 \int_{D(\rho,r)}|g_T(s)|^2\,dA(s)\geq\pi r^2.
\]
Summing these nonnegative integrals yields
\begin{equation}\label{30005276:detection}
 Z\ll r^{-2}(\log T)\mathcal E(\sigma-r,T)
   \ll_{\sigma,r,\eta}T^{2(1-\sigma)+2r+\eta}\log T.
\end{equation}
Take $\eta=\varepsilon/4$ and absorb the logarithm into
$T^{\varepsilon/4}$ for large $T$.  Dyadic summation bounds the positive
ordinates up to any height, with the finitely many small ordinates
absorbed in the constant; complex conjugation handles negative ones.
At $\sigma=1/2$, the standard total zero bound $O(T\log T)$ suffices.
At $\sigma=1$, the stipulated nontrivial zeros have $\beta<1$, so the
count is empty.  All constants may depend on $\sigma,\varepsilon$, as
allowed in the original statement. \hfill$\square$

This proof accounts for multiplicity through the local count, rather
than pretending that one disk detects a multiplicity proportional
amount of area.  It also avoids evaluating a function only on a
measure-zero set without a justified neighborhood estimate.

\paragraph{Attempt II: why the required exponent looks plausible.}
For $\Re s>1$, absolute convergence permits the exact expansion
\begin{equation}\label{30005276:coefficients}
 \zeta(s)M_T(s)-1=\sum_{n>X}a_X(n)n^{-s},\qquad
 X=\lfloor T\rfloor,\quad
 a_X(n)=\sum_{\substack{d\mid n\\d\leq X}}\mu(d).
\end{equation}
For $2\leq n\leq X$, the omitted coefficient is zero by
$\sum_{d\mid n}\mu(d)=0$; the constant coefficient has been subtracted.
Also $|a_X(n)|\leq\tau(n)$.  For fixed $u>1/2$ and sufficiently small
$\eta>0$, the elementary divisor bound $\tau(n)\ll_\eta n^{\eta/2}$ gives
\begin{equation}\label{30005276:diagonal}
 \sum_{n>X}|a_X(n)|^2n^{-2u}
 \ll_{u,\eta}X^{1-2u+\eta}.
\end{equation}
If a mean-square theorem with time length $T$ retained only this
coefficient-square term, it would give
$T X^{1-2u+\eta}=T^{2-2u+\eta}$, exactly the exponent needed by
\eqref{30005276:missing}.  Integrating over $u$ would not cost a new power
when $u\geq u_0>1/2$.

This is a motivation, not an application of a mean-square theorem.
Equation \eqref{30005276:coefficients} was proved only for $\Re s>1$.
Square summability of its coefficients with a weight does not establish
ordinary convergence of that Dirichlet series at the desired points,
nor does it justify a mean-square approximation uniform in $X=T$.
For fixed $X$, its summatory coefficients have main term
\[
 \sum_{n\leq Y}\sum_{\substack{d\mid n\\d\leq X}}\mu(d)
   =Y\sum_{d\leq X}\frac{\mu(d)}d+O(X),
\]
so absolute convergence across the line $1$ certainly has not followed
from the cancellation of the first $X$ coefficients.  Subtracting the
single constant coefficient does not change this issue.

\paragraph{Attempt III: an explicit off-diagonal obstruction.}
Even for a finite polynomial $D(t)=\sum_{L<n\leq2L}b_n n^{-it}$,
expanding the integral gives off-diagonal factors
\[
 \int_T^{2T}e^{it\log(m/n)}\,dt,
 \qquad
 \left|\int_T^{2T}e^{it\log(m/n)}\,dt\right|
 \leq\min\{T,2/|\log(m/n)|\}.
\]
Distinct integers do not make this integral vanish.  In particular,
adjacent integers near $L\gg T$ can be coherent throughout a substantial
part of the averaging interval.

A direct counterexample to a uniform ``diagonal only'' bound is
$b_n=n^{iT}$.  For $t=T+v$ with $0\leq v\leq1/4$,
\[
 e^{iv\log L}D(T+v)=\sum_{L<n\leq2L}e^{-iv\log(n/L)}.
\]
Every summand has real part at least $\cos((\log2)/4)>0$.  Hence
\[
 \int_T^{2T}|D(t)|^2\,dt\gg L^2,
 \qquad T\sum_{L<n\leq2L}|b_n|^2=TL.
\]
Taking $L/T\to\infty$ disproves a bound by an absolute constant times
the second expression.  This does \emph{not} disprove
\eqref{30005276:missing}: the coefficients here are not the arithmetic
coefficients of the mollifier.  It proves that the desired estimate
would require their arithmetic structure or a different cancellation
mechanism, rather than a general orthogonality principle.

\paragraph{Stress test.}
The attempt separates three levels: the area lemma is proved;
\eqref{30005276:diagonal} is a valid estimate of a formal diagonal expression;
and the passage from that expression to \eqref{30005276:area} is unproved.
An interchange of an infinite sum with an integral, or a passage from
fixed mollifier length to length $T$, is not asserted.  The loss
$2r$ in \eqref{30005276:detection} is made arbitrarily small only after
$\sigma,\varepsilon$ are fixed.  No uniformity as $\sigma\downarrow1/2$
is needed or claimed.

The companion computation checks the coefficient cancellation in
\eqref{30005276:coefficients} for finite $X$ and records exact finite
coefficient-square sums at rational exponents where possible.  It also
checks the coherent-polynomial construction numerically, explicitly as
a numerical check of an already proved elementary bound.  No computed
list of zeta zeros is promoted to a density theorem.  Even the full
density hypothesis would permit isolated zeros off the critical line;
this argument is not a proposed proof of the Riemann hypothesis.

\paragraph{Outcome.}
\textbf{Outcome: Conditional reduction.}

The density hypothesis has been reduced to the explicit area estimate
\eqref{30005276:missing} for a growing mollifier.  The exponent is explained
by a valid coefficient calculation, and the tempting invalid
mean-square step is falsified for general coefficients.  A proof of the
arithmetic off-diagonal estimate, with the required uniformity in the
mollifier length, is still missing.  No improved unconditional zero
density exponent or historical novelty claim is made.
% END RESEARCH ADDITION 128
\subsection{Sources}
\begin{itemize}\raggedright
\item[\textnormal{[S1]}]\hypertarget{source:30005276:1}{} Terence Tao, "A computation-outsourced discussion of zero density theorems for the Riemann zeta function", What's new, 7 July 2024.
\url{https://terrytao.wordpress.com/tag/riemann-zeta-function/}.
{\small\textit{Catalogue formulation source.}}
% BEGIN RESEARCH REFERENCES 128
\item[\textnormal{[E1]}]\hypertarget{extra:30005276:1}{} Larry Guth and James Maynard.
\emph{New large value estimates for Dirichlet polynomials}, arXiv:2405.20552v2, 7 April 2026. Theorem 1.2 and Section 13.1, including the local zero count and zero-detection reduction.
\url{https://arxiv.org/html/2405.20552v2}.
{\small\textit{Consulted 27 September 2026 for the indicated result or scope.}}
% END RESEARCH REFERENCES 128
\end{itemize}


\end{document}
